% Generated draft: requires mathematical and source review.
\documentclass[11pt]{article}
\usepackage[margin=25mm]{geometry}
\usepackage{fontspec}
\setmainfont{Latin Modern Roman}
\usepackage{amsmath,amssymb,mathtools,unicode-math}
\setmathfont{Latin Modern Math}
\usepackage{enumitem,needspace,xurl,hyperref}
\setlength{\parindent}{0pt}
\setlength{\parskip}{5pt}
\setlength{\emergencystretch}{3em}
\newcommand{\N}{\mathbb{N}}
\newcommand{\Z}{\mathbb{Z}}
\newcommand{\Q}{\mathbb{Q}}
\newcommand{\R}{\mathbb{R}}
\newcommand{\C}{\mathbb{C}}
\newcommand{\F}{\mathbb{F}}
\newcommand{\A}{\mathbb{A}}
\newcommand{\PP}{\mathbb{P}}
\newcommand{\E}{\mathbb{E}}
\newcommand{\eps}{\varepsilon}
\newcommand{\dd}{\,\mathrm{d}}
\begin{document}
\section*{UnsolvedMath ID: 800873}
\textbf{Source problem code:} catalogue entry
\subsection{Definitions and mathematical statement}
The supplied question assumes complex dimension \(n=2\) and \(0\leq\alpha<1\), but does not explicitly specify an ambient domain or define maximality. Under the usual local pluripotential interpretation, let \(\Omega\subset\C^2\) be a domain and let \(\varphi\colon\Omega\to\R\) be a plurisubharmonic function of class \(\mathcal C^{2,\alpha}\). Here plurisubharmonic means that the restriction to each complex line is subharmonic wherever defined. Interpret maximality by the comparison property: for every relatively compact open set \(G\Subset\Omega\) and every plurisubharmonic function \(v\) on \(G\),
\[
\left(\forall\zeta\in\partial G,\quad \limsup_{G\ni z\to\zeta}v(z)\leq\varphi(\zeta)\right)
\quad\Longrightarrow\quad
v\leq\varphi\text{ on }G.
\]
This domain and this interpretation of maximality require confirmation from the source context; no boundary regularity or boundary values are supplied.

Choose the normalization \(dd^c=i\partial\bar\partial\), which does not affect its support, and put
\[
S=\operatorname{supp}_{\Omega}(dd^c\varphi)
=\Omega\setminus\bigcup\{U\subset\Omega: U\text{ is open and }dd^c\varphi|_U=0\}.
\]
Let \(\mathbb D=\{\zeta\in\C:|\zeta|<1\}\). An embedded holomorphic disc in \(S\) is the image of a holomorphic embedding \(f\colon\mathbb D\to\Omega\) with \(f(\mathbb D)\subset S\). The assertion that \(\varphi\) is harmonic along this disc means
\[
\frac{\partial^2(\varphi\circ f)}{\partial\zeta\,\partial\bar\zeta}=0
\quad\text{on }\mathbb D.
\]
\textbf{Question (Dujardin; conditional local formulation).} For every such \(\alpha\), \(\Omega\), and maximal \(\varphi\), is \(S\) foliated, or laminated, by holomorphic discs along which \(\varphi\) is harmonic?

The source leaves the alternative ``foliated (or laminated)'' unresolved: it does not specify the required local product structure, transverse regularity, or whether singularities or exceptional points are allowed. Consequently, the intended geometric conclusion cannot be made uniquely precise from this record alone; mere coverage by holomorphic discs should not be substituted for that conclusion.
\subsection{Short English statement}
Dujardin asks whether, in complex dimension two, the support of \(dd^c\varphi\) for a maximal \(\mathcal C^{2,\alpha}\) function, \(0\leq\alpha<1\), is foliated or laminated by holomorphic discs on which \(\varphi\) is harmonic. The ambient setting and the precise foliation or lamination requirements need confirmation.
\subsection{Sources}
\begin{itemize}
\item Local UnsolvedMath registry: permanent ID \(800873\), exact source problem code \texttt{catalogue entry}, title \emph{Open Problems in Pluripotential Theory --- Question 4}. The supplied registry labels the problem open; this label is not an independently verified mathematical or literature assessment.
\item Dinew, Guedj, and Zeriahi, \emph{Open Problems in Pluripotential Theory} (2015), Question 4, attributed in the supplied statement to Dujardin. Supplied source URL: \url{https://arxiv.org/abs/1511.00705}.
\item Supplied provenance: extraction method \texttt{arxiv\_latex}; source-access date \texttt{2026-10-16}. This date and the source attribution are supplied provenance and have not been independently verified here.
\end{itemize}
\end{document}
